\section{多元函数积分学}
\subsection{多元函数积分的概念与性质}
\begin{mydef}{二重积分}{two-integral}

\end{mydef}
\begin{mydef}{三重积分}{three-integral}

\end{mydef}
\begin{mydef}{第一类曲线积分}{first-curve-integral}

\end{mydef}
\begin{mydef}{第二类曲线积分}{second-curve-integral}
\end{mydef}
\begin{mydef}{第一类曲面积分}{first-surface-integral}


\end{mydef}
\begin{mydef}{第二类曲面积分}{second-surface-integral}
\end{mydef}
\begin{conclusion}{几种多元函数积分的比较}{integral-compare}
\end{conclusion}
\begin{conclusion}{多元函数积分的性质}{integral-property}
\end{conclusion}
\begin{conclusion}{两类曲线积分的关系}{curve-integral-compare}

\end{conclusion}
\begin{conclusion}{两类曲面积分的关系}{surface-integral-compare}
\end{conclusion}
\begin{conclusion}{第一二类曲线(面)积分比较}{integral-compare}

\end{conclusion}
\begin{conclusion}{一些基本结论}{integral-conclusion}

\end{conclusion}
\subsection{多元函数积分变换为定积分}
\begin{conclusion}{曲线积分化为定积分}{curve-integral-to-definite}

\end{conclusion}

\begin{conclusion}{二重积分化为累次积分}{two-integral-to-definite}

\end{conclusion}

\begin{conclusion}{三重积分化为累次积分}{three-integral-to-definite}
\end{conclusion}

\begin{conclusion}{第一类曲面积分化为二重积分}{first-curve-integral-to-two}

\end{conclusion}

\begin{conclusion}{第二类曲面积分化为二重积分}{second-curve-integral-to-two}
\end{conclusion}
\subsection{重积分的变量替换}
\begin{conclusion}{平移变换}{translation}

\end{conclusion}

\begin{conclusion}{二重积分换元法}{two-integral-change}

\end{conclusion}
\begin{conclusion}{二重积分的极坐标变换}{two-integral-polar}

\end{conclusion}

\begin{conclusion}{三重积分的柱坐标变换}{three-integral-polar}
\end{conclusion}

\begin{conclusion}{三重积分的球坐标变换}{three-integral-polar}
\end{conclusion}
\subsection{多元函数积分计算技巧}
\begin{conclusion}{选择积分顺序}{integral-order}

\end{conclusion}

\subsubsection{对称性与奇偶性}
\begin{conclusion}{二重积分的对称性与奇偶性}{two-integral-symmetry}
\end{conclusion}

\begin{conclusion}{三重积分的对称性与奇偶性}{three-integral-symmetry}
\end{conclusion}

\begin{conclusion}{第一类曲线,曲面积分的对称性与奇偶性}{first-curve-integral-symmetry}
\end{conclusion}

\begin{conclusion}{第二类曲线,曲面积分的对称性与奇偶性}{second-curve-integral-symmetry}
\end{conclusion}

\begin{conclusion}{分块积分法}{integral-block}

\end{conclusion}

\begin{tablebox}{选择变量替换}

\end{tablebox}

\begin{conclusion}{计算曲线、曲面积分时利用曲线或曲面方程}{integral-geometry}
\end{conclusion}

\begin{conclusion}{计算第二类曲面积分时选择投影方向}{projection-direction}

\end{conclusion}
\subsection{多元函数积分的各种公式}
\begin{mydef}{平面上的单连通区域与区域的正向边界}{single-region}

\end{mydef}

\begin{formula}{格林公式}{Green-formula}

\end{formula}
\begin{formula}{高斯公式}{Gauss-formula}

\end{formula}
\begin{formula}{斯托克斯公式}{Stokes-formula}
\end{formula}

\begin{conclusion}{格林公式、高斯公式、斯托克斯公式的关系}{Green-Gauss-Stokes}

\end{conclusion}

\begin{mydef}{向量场的通量与散度}{flux-divergence}
\end{mydef}
\begin{mydef}{向量场的环流量与散度}{flux-divergence}
\end{mydef}

\begin{conclusion}{应用格林公式计算曲线积分}{Green-curve}

\end{conclusion}
\begin{conclusion}{应用高斯公式计算曲面积分}{Gauss-surface}
\end{conclusion}

\begin{conclusion}{应用斯托克斯公式计算曲线积分}{Stokes-curve}
\end{conclusion}

\begin{mydef}{曲线积分与路径无关问题}{path-independent}

\end{mydef}

\begin{mydef}{微分式的原函数问题}{indefinite-integral}

\end{mydef}

\begin{conclusion}{\(\int_{L}P(x,y)dx+Q(x,y)dy\)与路径无关时的特征}{path-independent-feature}

\end{conclusion}
\begin{conclusion}{如何判断曲线积分\(\int_{L}P(x,y)dx+Q(x,y)dy\)与路径无关}{path-independent-feature-judgement}

\end{conclusion}
\begin{conclusion}{积分与路径无关时求\(I = \int_{L}P(x,y)dx+Q(x,y)dy\)}{path-independent-integral}
\end{conclusion}
\begin{conclusion}{积分与路径无关时求原函数}{path-independent-indefinite}

\end{conclusion}
\subsection{多元函数积分的几何应用}
\begin{conclusion}{空间图形的体积}{space-volume}

\end{conclusion}
\begin{conclusion}{旋转面的面积}{rotation-area}

\end{conclusion}
\begin{conclusion}{一般曲面的面积}{surface-area}
\end{conclusion}
\subsection{多元函数积分的物理应用}
\begin{mydef}{质量、质心、形心与转动惯量}{mass-center-moment}

\end{mydef}
\begin{conclusion}{质量、质心与形心的计算}{mass-center}
\end{conclusion}

\begin{conclusion}{引力计算}{gravity-integral}

\end{conclusion}

\begin{conclusion}{功的计算}{work-integral}

\end{conclusion}
\begin{conclusion}{流量计算}{flow-integral}

\end{conclusion}
\subsection{其他技巧}